\documentclass[12pt,a4paper,reqno]{amsart}

% language
\usepackage[greek,english]{babel}
\usepackage[utf8]{inputenc}
\usepackage{alphabeta}

% change default names to greek
\addto\captionsenglish{
    \renewcommand{\contentsname}{Περιεχόμενα}
    \renewcommand{\refname}{Βιβλιογραφία}
    \renewcommand{\datename}{Ημερομηνία:}
    \renewcommand{\urladdrname}{Ιστοσελίδα}
}

% math 
\usepackage{amsmath,amsthm,amssymb,amscd}

% font
\usepackage[cal=euler]{mathalfa}
\usepackage{libertinus-type1}
% \usepackage{txfonts} % for upright greek letters
\usepackage{bm} % for bold symbols
\usepackage{bbm} % for the simply-looking bb symbols

% miscellaneous 
\usepackage{changepage} %for indenting environments
\usepackage{csquotes} % example: \textcquote{}
\usepackage{blkarray}
\setcounter{MaxMatrixCols}{20} % default for pmatrix is 10!!
\usepackage{ytableau}
\usepackage{array} %needed to increase the vertical length in a tabular

% drawing
\usepackage{tikz,tikz-cd}
\usetikzlibrary{shapes.misc, patterns, matrix, calc, intersections,positioning,arrows.meta}
\usepackage{graphics,graphicx}
\usepackage{float} % provides enhanced control and customization options for floating objects such as figures and tables

% colors
\usepackage{xcolor}
\definecolor{darkcandyapplered}{rgb}{0.64, 0.0, 0.0}
\definecolor{midnightblue}{rgb}{0.1, 0.1, 0.44}
\definecolor{mylightblue}{HTML}{336699}
\definecolor{burntorange}{rgb}{0.8, 0.33, 0.0}
\definecolor{iceberg}{rgb}{0.44, 0.65, 0.82}
\definecolor{applegreen}{rgb}{0.55, 0.71, 0.0}
\definecolor{canaryyellow}{rgb}{1.0, 0.94, 0.0}
\definecolor{lightgray}{rgb}{0.83, 0.83, 0.83}

% hrefs
\usepackage{hyperref}
\usepackage[noabbrev,capitalize]{cleveref}
\hypersetup{
    pdftoolbar=true,        
    pdfmenubar=true,        
    pdffitwindow=false,     
    pdfstartview={FitH},  % fits the width of the page to the window
    pdftitle={},
    pdfauthor={},
    pdfsubject={},
    pdfkeywords={},
    pdfnewwindow=true,  % links in new window
    colorlinks=true,  % false: boxed links; true: colored links
    linkcolor=darkcandyapplered,   % color of internal links
    citecolor=midnightblue,  % color of links to bibliography
    urlcolor=cyan,  % color of external links
    linktocpage=true  % changes the links from the section body to the page number
    }

% geometry
\textwidth=16cm 
\textheight=21cm 
\hoffset=-55pt 
\footskip=25pt

% thm envs (you might need to change the path)
\theoremstyle{definition}
\newtheorem{theorem}{Θεώρημα}
\newtheorem{proposition}[theorem]{Πρόταση}
\newtheorem{lemma}[theorem]{Λήμμα}
\newtheorem{corollary}[theorem]{Πόρισμα}
\newtheorem{conjecture}[theorem]{Εικασία}
\newtheorem{problem}[theorem]{Πρόβλημα}
\newtheorem*{claim}{Ισχυρισμός}
\newtheorem{observation}[theorem]{Παρατήρηση}
\newtheorem{definition}[theorem]{Ορισμός}
\newtheorem{question}[theorem]{Ερώτημα}
\newtheorem*{questions}{Ερωτήματα}
\newtheorem*{que}{Ερώτημα}
\newtheorem{example}[theorem]{Παράδειγμα}
\newtheorem{exercise}{Άσκηση}
\newtheorem*{zorn}{Λήμμα του Zorn}
\newtheorem*{example_6.6}{Παράδειγμα~6.6}

\theoremstyle{remark}
\newtheorem*{remark}{Παρατήρηση}

% fixes the correct numbering of environments
\numberwithin{theorem}{section}
\numberwithin{exercise}{section}
\numberwithin{equation}{section}

% math ops 
% fields
\newcommand{\NN}{\mathbbmss{N}} 
\newcommand{\ZZ}{\mathbbmss{Z}} 
\newcommand{\QQ}{\mathbbmss{Q}} 
\newcommand{\RR}{\mathbbmss{R}} 
\newcommand{\CC}{\mathbbmss{C}} 
\newcommand{\KK}{\mathbbmss{K}} 
\newcommand{\FF}{\mathbbmss{F}} 
\newcommand{\HH}{\mathbbmss{H}} 

% symmetric group
\newcommand{\fS}{\mathfrak{S}}  

% calligraphic 
\newcommand{\aA}{\mathcal{A}} 
\newcommand{\bB}{\mathcal{B}}
\newcommand{\cC}{\mathcal{C}}
\newcommand{\dD}{\mathcal{D}}
\newcommand{\eE}{\mathcal{E}}
\newcommand{\fF}{\mathcal{F}}
\newcommand{\hH}{\mathcal{H}}
\newcommand{\iI}{\mathcal{I}}
\newcommand{\lL}{\mathcal{L}}
\newcommand{\oO}{\mathcal{O}}
\newcommand{\pP}{\mathcal{P}}
\newcommand{\qQ}{\mathcal{Q}}
\newcommand{\sS}{\mathcal{S}}
\newcommand{\mM}{\mathcal{M}}
\newcommand{\uU}{\mathcal{U}}

% bold
\newcommand{\bfa}{\mathbf{a}}
\newcommand{\bfe}{\mathbf{e}}
\newcommand{\bfF}{\pmb{F}}
\newcommand{\bfR}{\pmb{R}}
\newcommand{\bfv}{\mathbf{v}}
%\newcommand{\bfx}{\bm{x}}
%\newcommand{\bfx}{\mathbf{x}} 
\newcommand{\bfx}{\pmb{x}}
\newcommand{\bfX}{\pmb{X}}
\newcommand{\bfy}{\pmb{y}}
\newcommand{\bfz}{\pmb{z}}

% roman
\newcommand{\rmA}{\mathrm{A}}
\newcommand{\rmB}{\mathrm{B}}
\newcommand{\rmC}{\mathrm{C}}
\newcommand{\rmD}{\mathrm{D}} 
\newcommand{\rmF}{\mathrm{F}} 
\newcommand{\rmH}{\mathrm{H}} 
\newcommand{\rmh}{\mathrm{h}} 
\newcommand{\rmI}{\mathrm{I}} 
\newcommand{\rmK}{\mathrm{K}}
\newcommand{\rmM}{\mathrm{M}}
\newcommand{\rmP}{\mathrm{P}}  
\newcommand{\rmp}{\mathrm{p}}  
\newcommand{\rmQ}{\mathrm{Q}}  
\newcommand{\rmR}{\mathrm{R}}
\newcommand{\rmS}{\mathrm{S}}
\newcommand{\rmT}{\mathrm{T}}
\newcommand{\rmU}{\mathrm{U}}
\newcommand{\rmV}{\mathrm{V}}
\newcommand{\rmY}{\mathrm{Y}}
\newcommand{\rmZ}{\mathrm{Z}}
\newcommand{\rmz}{\mathrm{z}}

% arrows and symbols 
\renewcommand{\to}{\rightarrow}
\newcommand{\toto}{\longrightarrow}
\newcommand{\mapstoto}{\longmapsto}
\newcommand{\then}{\Rightarrow}
\newcommand{\IFF}{\Leftrightarrow}
\newcommand{\tl}{\tilde}
\newcommand{\wtl}{\widetilde}
\newcommand{\ol}{\overline}
\newcommand{\ul}{\underline}
\newcommand{\oldemptyset}{\emptyset}
\renewcommand{\emptyset}{\varnothing}
\DeclareMathSymbol{\Arg}{\mathbin}{AMSa}{"39} % for arguments 
\newcommand{\onto}{\ensuremath{\twoheadrightarrow}}
\newcommand{\tle}{\trianglelefteq}
\newcommand{\tge}{\trianglerighteq}

% custom ops
\newcommand{\rmKer}{\mathrm{Ker}}
\newcommand{\rmIm}{\mathrm{Im}}
\newcommand{\lcm}{\mathrm{lcm}}
\newcommand{\GL}{\mathrm{GL}}
\newcommand{\ev}{\mathrm{ev}}
\newcommand{\End}{\mathrm{End}}
\newcommand{\rmid}{\mathrm{id}}
\newcommand{\rmspan}{\mathrm{span}}
\newcommand{\Hom}{\mathrm{Hom}}
\newcommand{\Ann}{\mathrm{Ann}}
\newcommand{\Set}{\pmb{\mathrm{Set}}}
\newcommand{\Rmod}{\pmb{\mathrm{mod}}}
\newcommand{\rmrank}{\mathrm{rank}}
\newcommand{\mSpec}{\mathrm{mSpec}}
\newcommand{\Spec}{\mathrm{Spec}}
\newcommand{\tor}{\mathrm{tor}}

% absolute value symbol
\usepackage{mathtools} 
\DeclarePairedDelimiter\abs{\lvert}{\rvert}%
\DeclarePairedDelimiter\norm{\lVert}{\rVert}%
\makeatletter
\let\oldabs\abs
\def\abs{\@ifstar{\oldabs}{\oldabs*}}

% tensor symbol
\newcommand{\tensor}[1]{%
  \mathbin{\mathop{\otimes}\limits_{#1}}%
}

% permutation cycle notation
\ExplSyntaxOn
\NewDocumentCommand{\cycle}{ O{\;} m }
 {
  (
  \alec_cycle:nn { #1 } { #2 }
  )
 }

\seq_new:N \l_alec_cycle_seq
\cs_new_protected:Npn \alec_cycle:nn #1 #2
 {
  \seq_set_split:Nnn \l_alec_cycle_seq { , } { #2 }
  \seq_use:Nn \l_alec_cycle_seq { #1 }
 }
\ExplSyntaxOff

% setminus symbol
\newcommand{\mysetminusD}{\hbox{\tikz{\draw[line width=0.6pt,line cap=round] (3pt,0) -- (0,6pt);}}}
\newcommand{\mysetminusT}{\mysetminusD}
\newcommand{\mysetminusS}{\hbox{\tikz{\draw[line width=0.45pt,line cap=round] (2pt,0) -- (0,4pt);}}}
\newcommand{\mysetminusSS}{\hbox{\tikz{\draw[line width=0.4pt,line cap=round] (1.5pt,0) -- (0,3pt);}}}
\newcommand{\sm}{\mathbin{\mathchoice{\mysetminusD}{\mysetminusT}{\mysetminusS}{\mysetminusSS}}}

% miscellaneous commands
\newcommand{\defn}[1]{{\color{mylightblue}{#1}}}
\newcommand{\toDo}{{\bf\color{red} TODO}}
\newcommand{\toCite}{{\bf\color{green} CITE}}
\newcommand*{\vertbar}{\rule[-1ex]{0.5pt}{2.5ex}} % for matrices with column vectors
\newcommand*{\horzbar}{\rule[.5ex]{2.5ex}{0.5pt}} % for matrices with row vectors
\newcommand{\myblue}[1]{{\color{iceberg}{#1}}}
\newcommand{\myorange}[1]{{\color{burntorange}{#1}}}
\newcommand{\mygreen}[1]{{\color{applegreen}{#1}}}
\newcommand{\myred}[1]{{\color{darkcandyapplered}{#1}}}

% ferrer's diagram
\newcommand{\fdiagram}[1]{
    \begin{tikzpicture}[scale=.7]
        \fill foreach \Z [count=\Y] in {#1}
        {foreach \X in {1,...,\Z} 
        {(\X,-\Y) circle[radius=3pt]}};
    \end{tikzpicture}
}

%
\usepackage{pgfplots}
\pgfplotsset{compat=1.18}

%
\makeatletter
\def\mathcolor#1#{\@mathcolor{#1}}
\def\@mathcolor#1#2#3{%
  \protect\leavevmode
  \begingroup
    \color#1{#2}#3%
  \endgroup
}
\makeatother

% restriction
\newcommand\restr[2]{{% we make the whole thing an ordinary symbol
  \left.\kern-\nulldelimiterspace % automatically resize the bar with \right
  #1 % the function
  \littletaller % pretend it's a little taller at normal size
  \right|_{#2} % this is the delimiter
  }}
\newcommand{\littletaller}{\mathchoice{\vphantom{\big|}}{}{}{}}

% 
\newenvironment{nouppercase}{%
  \let\uppercase\relax%
  \renewcommand{\uppercasenonmath}[1]{}}{}

% titlepage
\title{226: Θεωρία Δακτυλίων και Modules}
\author[Β.~Δ. Μουστακας]{Βασίλης Διονύσης Μουστάκας \\ Πανεπιστήμιο Κρήτης}
\date{28 Απριλίου 2026}
% \urladdr{\href{https://sites.google.com/view/vasmous}{https://sites.google.com/view/vasmous}}

\begin{document}

\begingroup
\def\uppercasenonmath#1{} % this disables uppercase title
\let\MakeUppercase\relax % this disables uppercase authors
\maketitle
\endgroup

\setcounter{section}{7}
\setcounter{theorem}{7}
\setcounter{equation}{5}
\begin{center}
    \textbf{7. Πρότυπα πάνω από περιοχές κύριων ιδεωδών
} (Συνέχεια)
\end{center}

\begin{proof}[Συνέχεια απόδειξης του Θεωρήματος 7.3 (Μοναδικότητα αναλλοίωτων παραγόντων)]
θα δείξουμε ότι $(a_i) = (b_i)$ για κάθε $1 \le i \le \ell$. Σταθεροποιούμε ένα υποδείκτη $1 \le i \le \ell$ και παρατηρούμε ότι 
    \[
    a_iM_1 = a_iM_2 = \cdots = a_iM_i = 0
    \]
    Πράγματι, εξ υποθέσεως 
    \[
    (a_i) \subseteq (a_{i-1}) \subseteq \cdots \subseteq (a_1).
    \]
    Όμως, για κάθε $1 \le j \le i$, $a_jM_j = 0$, διότι το $a_j$ είναι γεννήτορας του $\Ann_R(M_j)$ και το ζητούμενο έπεται. Επομένως, 
    \[
    a_iM \cong a_iM_{i+1} \oplus \cdots \oplus a_iM_\ell
    \]
    και εξ' υποθέσεως
    \begin{equation}  
        \label{eq:structure_help}
        a_iM \cong a_iN_1 \oplus a_iN_2 \oplus \cdots \oplus a_iN_\ell.
    \end{equation}
    Επειδή τα $a_iM_j$ και $a_iN_j$ είναι κυκλικά πρότυπα, απαριθμώντας γεννήτορες και στα δύο μέλη της Ταυτότητας \eqref{eq:structure_help}, έπεται ότι τουλάχιστον $i$ προσθεταίοι του δεξιού μέλους είναι μηδενικά πρότυπα. 
    
    Ισχυριζόμαστε ότι $a_iN_i = 0$. Πράγματι, αν $a_iN_i \neq 0$, τότε $a_i \notin \Ann_R(N_i) = (b_i)$. Όμως, εξ υποθέσεως  
    \[
    (b_\ell) \subseteq (b_{\ell-1}) \subseteq \cdots \subseteq (b_{i+1}) \subseteq (b_i)
    \]
    και γι αυτό $a_i \notin (b_j) = \Ann_R(N_j)$ για κάθε $j \ge i$. Επομένως, $a_iN_j \neq 0$, για κάθε $j \ge i$. Με άλλα λόγια, βρήκαμε τουλάχιστον $\ell -i +1$ μη μηδενικούς προσθεταίους. Ισοδύναμα, αυτό σημαίνει ότι υπάρχουν το πολύ $i-1$ μηδενικοί προσθεταίοι, το οποίο είναι αδύνατο. 
    
    Επειδή $a_iN_i = 0$ έπεται ότι $(a_i) \subseteq (b_i)$. Ομοίως, $(b_i) \subseteq (a_i)$. Άρα $(a_i) = (b_i)$ και το ζητούμενο έπεται.
\end{proof}

Μένει να αποδείξουμε την ύπαρξη και την μοναδικότητα της διάσπασης σε στοιχειώδεις διαιρέτες. Τώρα που γνωρίζουμε την διάσπαση σε αναλλοίωτους παράγοντες, το πρώτο είναι ουσιαστικά εφαρμογή του κινεζικού θεωρήματος υπολοίπων.

\begin{proof}[Απόδειξη του Θεωρήματος 7.3 (Ύπαρξη στοιχειωδών διαιρετών)]
    Υποθέτουμε ότι 
    \[
    M \cong R^n \oplus R/(a_1) \oplus R/(a_2) \oplus \cdots \oplus R/(a_\ell),
    \]
για κάποια $n, \ell \ge 0$ και κάποια μη μηδενικά, μη αντιστρέψιμα $a_1, a_2, \dots, a_\ell \in R$ τέτοια ώστε $a_i \mid a_{i+1}$, για κάθε $1 \le i \le \ell-1$. Επειδή ο $R$ είναι περιοχή κύριων ιδεωδών, από το Θεώρημα 2.12 έπεται ότι είναι και περιοχή μονοσήμαντης παραγοντοποίησης και γι αυτό 
\[
a_i = \prod_{\substack{p \in R \\ \text{το p είναι πρώτο στοιχείο του $R$}}} p^{\lambda_i^{(p)}}.
\]
Το (πεπερασμένο) πλήθος των διαφορετικών πρώτων στοιχείων που εμφανίζονται είναι τουλά\-χιστον $i$, διότι $a_1 \mid a_2 \mid \cdots \mid a_\ell$. Από το κινέζικο θεώρημα υπολοίπων, έπεται ότι 
\begin{equation}
    \label{eq:elementary_divisors_help}
    R/(a_i) \cong \bigoplus_{\substack{p \in R \\ \text{το p είναι πρώτο στοιχείο του $R$}}} R/(p^{\lambda_i^{(p)}})
\end{equation}
για κάθε $1 \le i \le \ell$. Η διάσπαση (7.2) έπεται ομαδοποιώντας τους προσθεταίους στην διάσπαση (7.1) ως προς πρώτα στοιχεία, χρησιμοποιώντας την Ταυτότητα \eqref{eq:elementary_divisors_help}.
\end{proof}

Η διαδικασία αυτή μπορεί να \textquote{αντιστραφεί} για να κατασκευάσει κανείς τους αναλλοίωτους παράγοντες δοθέντων των στοιχειωδών διαιρετών, όπως έγινε στο Παράδειγμα 7.4. Έτσι έπεται η μοναδικότητα της διάσπασης σε στοιχειώδεις διαιρέτες. Οι λεπτομέρειες αφήνονται ως άσκηση.

\begin{remark}
    Η ακολουθία θετικών ακεραίων $\lambda = (\lambda_1,\lambda_2, \dots, \lambda_k)$ που εμφανίζεται στη διάσπαση σε στοιχειώδεις διαιρέτες και ικανοποιεί τη συνθήκη 
    \[
    \lambda_1 \ge \lambda_2 \ge \cdots \ge \lambda_k
    \]
    ονομάζεται \defn{διαμέριση ακεραίου} (integer partition). Αν $\lambda_1 + \lambda_2 + \cdots + \lambda_k = n$, τότε λέμε ότι το $\lambda$ είναι διαμέριση του $n$. Για παράδειγμα, οι διαμερίσεις του 4 είναι 
    \[
    (1,1,1,1), \ (2,1,1), \ (2,2), \ (3,1), \ (4).
    \]
    Οι διαμερίσεις ακεραίων παίζουν κεντρικό ρόλο στη θεωρία αριθμών, την αλγεβρική συνδυαστική και τη θεωρία αναπαραστάσεων των συμμετρικών ομάδων (βλ., για παράδειγμα, \cite[Ενότητα~6]{StaAC}).
\end{remark}

Το Θεώρημα 7.3 έχει δύο πολύ σημαντικές εφαρμογές. Στην περίπτωση όπου $R=\ZZ$, ταξινομεί όλες τις πεπερασμένα παραγόμενες αβελιανές ομάδες. 
\begin{example}
\label{ex:abelian_groups}
Ας ταξινομήσουμε όλες τις αβελιανές ομάδες τάξης $1800 = 2^33^25^2$ χρησιμοποιώντας τους στοιχειώδεις διαιρέτες. Έχουμε 
\[
\renewcommand{\arraystretch}{1.2}
\begin{array}{c|c|c}
    p^n & \text{διαμερίσεις του $n$} & \text{αβελιανές ομάδες τάξης $p^n$} \\ \hline
    2^3 & (1,1,1)                    & \ZZ_2\oplus\ZZ_2\oplus\ZZ_2 \\ 
        & (2,1)                      & \ZZ_4 \oplus \ZZ_2 \\
        & (3)                        & \ZZ_8 \\\hline
    3^2 & (1,1)                      & \ZZ_3 \oplus \ZZ_3 \\
        & (2)                        & \ZZ_9 \\ \hline 
    5^2 & (1,1)                      & \ZZ_5 \oplus \ZZ_5 \\ 
        & (2)                        & \ZZ_{25}
\end{array} \ .
\]
Παίρνοντας όλους τους πιθανούς συνδυασμούς των αβελιανών ομάδων από τη δεξιά στήλη προκύπτουν όλες οι αβελιανές ομάδες τάξης $1800$
\[
\renewcommand{\arraystretch}{1.2}
\begin{array}{c|c|c|c}
    \text{αβελιανή ομάδα} & \lambda^{(2)} & \lambda^{(3)} & \lambda^{(5)} \\ \hline
    \ZZ_2 \oplus \ZZ_2 \oplus \ZZ_2 \oplus \ZZ_3 \oplus \ZZ_3 \oplus \ZZ_5 \oplus \ZZ_5 & (1,1,1) & (1,1) & (1,1) \\ \hline
    \ZZ_2 \oplus \ZZ_2 \oplus \ZZ_2 \oplus \ZZ_3 \oplus \ZZ_3 \oplus \ZZ_{25} & (1,1,1) & (1,1) & (2) \\ \hline
    \ZZ_2 \oplus \ZZ_2 \oplus \ZZ_2 \oplus \ZZ_9 \oplus \ZZ_5 \oplus \ZZ_5 & (1,1,1) & (2) & (1,1) \\ \hline
    \ZZ_2 \oplus \ZZ_2 \oplus \ZZ_2 \oplus \ZZ_9 \oplus \ZZ_{25} & (1,1,1) & (2) & (2) \\ \hline
    \ZZ_4 \oplus \ZZ_2 \oplus \ZZ_3 \oplus \ZZ_3 \oplus \ZZ_5 \oplus \ZZ_5 & (2,1) & (1,1) & (1,1) \\ \hline
    \ZZ_4 \oplus \ZZ_2 \oplus \ZZ_3 \oplus \ZZ_3 \oplus \ZZ_{25} & (2,1) & (1,1) & (2) \\ \hline
    \ZZ_4 \oplus \ZZ_2 \oplus \ZZ_9 \oplus \ZZ_5 \oplus \ZZ_5 & (2,1) & (2) & (1,1) \\ \hline
    \ZZ_4 \oplus \ZZ_2 \oplus \ZZ_9 \oplus \ZZ_{25} & (2,1) & (2) & (2) \\ \hline
    \ZZ_8 \oplus \ZZ_3 \oplus \ZZ_3 \oplus \ZZ_5 \oplus \ZZ_5 & (3) & (1,1) & (1,1) \\ \hline
    \ZZ_8 \oplus \ZZ_3 \oplus \ZZ_3 \oplus \ZZ_{25} & (3) & (1,1) & (2) \\ \hline
    \ZZ_8 \oplus \ZZ_9 \oplus \ZZ_5 \oplus \ZZ_5 & (3) & (2) & (1,1) \\ \hline
    \ZZ_8 \oplus \ZZ_9 \oplus \ZZ_{25} & (3) & (2) & (2)
\end{array} \ .
\]
\end{example}

Στην περίπτωση όπου $R = F[x]$, όπου $F$ είναι ένα σώμα, το Θεώρημα 7.3 ταξινομεί όλα τα πεπερασμένα παραγόμενα $F[x]$-πρότυπα, δηλαδή διανυσματικούς χώρους πάνω από το $F$ εφοδιασμένους με κάποιο $T \in \End_F(V)$. Στην περίπτωση αυτή, η διάσπαση σε
\begin{itemize}
    \item αναλλοίωτους παράγοντες ονομάζεται \emph{ρητή κανονική μορφή} (rational canonical form)
    \item στοιχειώδεις διαιρέτες ονομάζεται \emph{κανονική μορφή Jordan} (Jordan canonical form).
\end{itemize} 
Οι δύο αυτές διασπάσεις είναι πολύ σημαντικές στην γραμμική άλγεβρα. Στις σημειώσεις αυτές δε θα τις μελετήσουμε και γι αυτό παραπέμπουμε στα \cite[Ενότητα 3, Μέρος II]{EKPA_LA}, \cite[Κεφάλαιο 7]{EKPA_PID} και \cite[Κεφάλαιο 12]{DF04}.


\newpage 
\setcounter{section}{8}
\setcounter{theorem}{0}
\setcounter{equation}{0}
\begin{center}
    \textbf{8. Δακτύλιοι και πρότυπα της Noether
} 
\end{center}

Από εδώ και στο εξής υποθέτουμε ότι κάθε δακτύλιος είναι μεταθετικός.

Στις προηγούμενες ενότητες είδαμε ότι σε μία περιοχή κύριων ιδεωδών
\begin{itemize}
    \item κάθε ιδεώδες είναι πεπερασμένα παραγόμενο, και πιο συγκεκριμένα κυκλικό (βλ. Ορισμό 2.5)
    \item κάθε αύξουσα ακολουθία ιδεωδών είναι τελικά σταθερή (βλ. Λήμμα 2.13)
    \item κάθε μη κενό σύνολο ιδεωδών έχει μεγιστικό στοιχείο (αποδεικνύεται με τρόπο παρόμοιο με το Λήμμα 5.8 (πως;)).
\end{itemize}
Η επόμενη πρόταση μας πληροφορεί ότι αυτές οι τρεις ιδιότητες είναι ισοδύναμες.

\begin{proposition}
    \label{prop:noetherian}
    Έστω $M$ ένα $R$-πρότυπο. Τα ακόλουθα είναι ισοδύναμα.
    \begin{itemize}
        \item[(1)] Κάθε αύξουσα ακολουθία υποπροτύπων του $M$
        \[
        M_1 \subseteq M_2 \subseteq \cdots 
        \]
        είναι τελικά σταθερή, δηλαδή υπάρχει θετικός ακέραιος $k$ τέτοιος ώστε 
        \[
        M_k = M_{k+1} = \cdots.
        \]
        \item[(2)] Κάθε μη κενό σύνολο υποπροτύπων του $M$ έχει μεγιστικό στοιχείο.
        \item[(3)] Κάθε υποπρότυπο του $M$ είναι πεπερασμένα παραγόμενο.
    \end{itemize}
    Το (1) ονομάζεται \defn{συνθήκη αύξουσας αλυσίδας} (ascending chain condition).
\end{proposition}
\begin{proof}[Απόδειξη]
    Για τo $(1) \then (2)$, θεωρούμε ένα μη κενό σύνολο $S$ υποπροτύπων του $M$ και έστω $M_1 \in S$. Αν το $M_1$ είναι μεγιστικό στο $S$ τελειώσαμε. Διαφορετικά, υπάρχει $M_2 \in S$ τέτοιο ώστε $M_1 \subseteq M_2$. Αν το $M_2$ είναι μεγιστικό, τότε τελειώσαμε. Διαφορετικά, υπάρχει $M_3 \in S$ τέτοιο ώστε $M_2 \subseteq M_3$ κοκ. Συνεχίζοντας έτσι προκύπτει μία αύξουσα ακολουθία υποπροτύπων του $M$ 
    \[
    M_1 \subseteq M_2 \subseteq M_3 \subseteq \cdots  
    \]
    Εξ' υποθέσεως, υπάρχει $k \ge 1$ τέτοιο ώστε 
    \[
    M_1 \subseteq M_2 \subseteq \cdots \subseteq M_k = M_{k+1} = \cdots. 
    \]
    Το $M_k$ είναι το ζητούμενο μεγιστικό στοιχείο του $S$.

    Για το $(2) \then (3)$, έστω $N$ ένα υποπρότυπο του $M$ και $S$ το σύνολο όλων των πεπερασμένα παραγόμενων υποπροτύπων του $M$ που περιέχονται στο $N$. Αυτό είναι μη κενό, διότι $0 \in S$. Εξ' υποθέσεως, το $S$ έχει κάποιο μεγιστικό στοιχείο $L$. Αν $N \neq L$, τότε υπάρχει κάποιο στοιχείο $a \in N \sm L$. Όμως, $L + (a) \in S$ (γιατί;) και γι αυτό $L = L + (a)$, πράγμα αδύνατο. Άρα, $N = L$ και κατά συνέπεια το $N$ είναι πεπερασμένα παραγόμενο.

    Για το $(3) \then (2)$, έστω 
    \[
    M_1 \subseteq M_2 \subseteq \cdots 
    \]
    μια αύξουσα ακολουθία υποπροτύπων του $M$. Θεωρούμε το υποπρότυπο 
    \[
    N = M_1 \cup M_2 \cup \cdots
    \]
    (βλ. και την απόδειξη του Λήμματος 2.13). Εξ' υποθέσεως, 
    \[
    N = R\{a_1, a_2, \dots, a_n\},
    \]
    όπου $a_j \in M_{i_j}$, για κάποια $i_j \ge 1$, για κάθε $1 \le j \le n$. Αν $k = \max\{i_j : 1\le j \le n\}$, τότε $a_j \in M_k$, για κάθε $1 \le j \le n$ και γι αυτό $N \subseteq M_k$. Επομένως, 
    \[
    M_1 \subseteq M_2 \subseteq \cdots \subseteq M_k = M_{k+1} = \cdots.
    \]
\end{proof} 

\begin{definition}
    \label{def:noetherian}
    Ένα $R$-πρότυπο που ικανοποιεί την συνθήκη αύξουσας αλυσίδας\footnote{Και κατ' επέκταση και τις υπόλοιπες δύο συνθήκες της Πρότασης \ref{prop:noetherian}.} ονομάζεται \defn{πρότυπο της Noether} (Noetherian module). Ειδικότερα, ο $R$ ονομάζεται \defn{δακτύλιος της Noether} αν το κανονικό $R$-πρότυπο είναι πρότυπο της Noether.
\end{definition}

\begin{example}
    \label{ex:noetherian}
    \leavevmode
    \begin{itemize}
        \item[(α)] Κάθε πεπερασμένο πρότυπο είναι πρότυπο της Noether.
        \item[(β)] Κάθε σώμα είναι δακτύλιος της Noether.
        \item[(γ)] Κάθε περιοχή κύριων ιδεωδών είναι δακτύλιος της Noether. Ειδικότερα, οι δακτύλιοι $\ZZ$, $\ZZ[i]$ και $F[x]$, για κάποιο σώμα $F$, είναι δακτύλιοι της Noether.
        \item[(δ)] Ο πολυωνυμικός δακτύλιος σε (αριθμήσιμα) άπειρες μεταβλητές $\ZZ[x_1, x_2, \dots]$ δεν είναι δακτύλιος της Noether. Πράγματι,
        \begin{itemize} 
            \item η ακολουθία ιδεωδών 
            \[
            (x_1) \subset (x_1, x_2) \subset \cdots 
            \]
            δεν σταθεροποιείται.
            \item το σύνολο ιδεωδών 
            \[
            \{(x_1), (x_1, x_2), \dots \}
            \]
            δεν έχει μεγιστικό στοιχείο, και 
            \item το ιδεώδες 
            \[
            (x_1, x_2, \dots)
            \]
            δεν είναι πεπερασμένα παραγόμενο.
        \end{itemize}
        (Γιατί;)
    \end{itemize}
\end{example}


\begin{thebibliography}{99}
    \bibitem{EKPA_LA}
    Δ. Α. Βάρσος, Δ. Ι. Δεριζιώτης, Ι. Π. Εμμανουήλ, Μ. Π. Μαλιάκας, Α. Δ. Μελάς και Ο. Π. Ταλέλλη, 
    \emph{Μια Εισαγωγή στην Γραμμική Άλγεβρα}, 
    Εκδόσεις Σοφία, 2012.
    %
    \bibitem{DF04}
    D. S. Dummit και R. M. Foote, 
    \textit{Abstract algebra},
    third edition, Wiley, 2004. 
    %
    \bibitem{EKPA_PID}
    Μ. Π. Μαλιάκας, και Ο. Π. Ταλέλλη, 
    \emph{Πρότυπα πάνω από Περιοχές Κυρίων Ιδεωδών και Εφαρμογές}, 
    Εκδόσεις Σοφία, 2009.
    %
    \bibitem{StaAC}
    R. P. Stanley,
    \emph{Algebraic Combinatorics; Walks, Trees, Tableaux and More},
    second edition, Undergraduate Texts in Mathematics, Springer, 2018.
\end{thebibliography}
\end{document}